% !TEX root = ../LAFROX_AnalisisFinanciero_2_20261113.tex
\chapter{Análisis económico-financiero}
\label{cap:e2}

El proyecto cuesta 210.840,39 UF y se ofrece a 246.910,17 UF netas. Deja un margen de 14,61\,\% y un VAN de 25.642,58 UF antes de impuestos. La caja financiada se mantiene positiva los 56 meses. Las cifras reflejan el modelo financiero, incluida la provisión de 384 UF para la instalación en Talca, absorbida por LafroX sin subir el precio. El capítulo separa las cifras del modelo de la asignación técnica del T-14, T-15, T-11 y T-19 (LafroX, 2026a, 2026c, 2026d, 2026f).

\parte{secciones/2-1_costos_frente_a_precio}
\parte{secciones/2-2_evaluacion_financiera}
\parte{secciones/2-3_curva_s}
\parte{secciones/2-4_estrategia_adquisiciones}
\parte{secciones/2-5_costos_operacionales}
\parte{secciones/2-6_mantencion_soporte}
\parte{secciones/2-7_valorizacion_innovaciones}

\section*{Referencias}
\phantomsection
\addcontentsline{toc}{section}{Referencias}

\noindent Distribuidora Puelche S.A. (2026a). \textit{Bases Administrativas para la preparación de la propuesta técnico-económica, licitación TFEP-01/2026}.

\noindent Distribuidora Puelche S.A. (2026b). \textit{Bases Técnicas Transversales, licitación TFEP-01/2026}.

\noindent LafroX. (2026a). \textit{Subdocumento 4: Arquitectura lógica y física de la solución, Anexo 4-W y Formulario T-11}.

\noindent LafroX. (2026b). \textit{Subdocumento 6: Metodologías, Anexo 6.B y Formulario T-9}.

\noindent LafroX. (2026c). \textit{Subdocumento 7: Plan de trabajo y Formularios T-14 y T-15}.

\noindent LafroX. (2026d). \textit{Subdocumento 13: Innovaciones y Formulario T-19}.

\noindent LafroX. (2026e). \textit{Subdocumento 8: Plan de riesgos, Anexo 8.H y Formulario T-16}.

\noindent LafroX. (2026f). \textit{Modelo financiero en planilla de cálculo, Entregable 3 de la Oferta Económica}.

\section*{Declaración de uso de IA}
\phantomsection
\addcontentsline{toc}{section}{Declaración de uso de IA}

Codex y Claude sirvieron como apoyo para ordenar la información y cotejar cifras. Antes de incorporar cada sección a este entregable, las personas responsables revisaron los datos con atención, compararon los resultados con las fuentes correspondientes y corrigieron cualquier diferencia o ambigüedad. La Tabla~\ref{tab:e2-ia} resume qué aspectos fue revisado por cada responsable (Distribuidora Puelche S.A., 2026a, art.~13.5).

\begin{tablalafrox}{Uso de IA por sección.}{tab:e2-ia}{F{0.09} F{0.11} F{0.28} F{0.09} F{0.13} F{0.30}}{\textbf{Sección} & \textbf{Herr.} & \textbf{Finalidad} & \textbf{Texto} & \textbf{Diagramas} & \textbf{Revisión humana}}
Introducción & Claude & Resumen general del capítulo & Medio & Ninguno & Tomás Pérez revisó las cifras principales y la relación entre la introducción y las secciones siguientes. Corrigió expresiones que podían resultar ambiguas y confirmó la coherencia con el modelo financiero. \\
2.1 & Codex y Claude & Costos, precios y márgenes & Alto & Ninguno & Alex Aravena verificó subtotales, conversiones de moneda y márgenes contra el modelo financiero. También revisó las partidas absorbidas por LafroX y el cumplimiento del tope de margen. \\
2.2 & Claude & Indicadores y flujos financieros & Alto & Ninguno & Tomás Pérez recalculó el VAN, la TIR y el punto de equilibrio, además de contrastar los flujos mensuales y anuales. Ajustó la explicación de la tasa de descuento y las sensibilidades. \\
2.3 & Codex y Claude & Flujo financiado y Curva S & Alto & Bajo & Alex Aravena comparó hitos, intereses y saldos con las hojas de financiamiento y Curva S. Revisó el mínimo de caja y comprobó que el gráfico no dependiera de la línea de crédito de respaldo. \\
2.4 & Codex y Claude & Adquisiciones y proveedores & Alto & Ninguno & Álvaro Catalán revisó montos, proveedores y cronograma contra el modelo financiero y los documentos T-11 y T-14. También aclaró qué equipos adquiría LafroX y cuáles correspondía instalar al cliente. \\
2.5 & Codex y Claude & Modelo de costos operacionales & Alto & Ninguno & Maximiliano Miño cotejó costos mensuales, recursos humanos, servicios y gastos administrativos con el modelo financiero. Revisó las unidades y corrigió posibles doble conteos. \\
2.6 & Codex y Claude & Soporte y mantención & Alto & Ninguno & Alex Aravena contrastó perfiles, dedicaciones y horas con el T-14 y el T-15. Verificó los costos de mantención y que no se repitieran partidas en el presupuesto de Operación. \\
2.7 & Codex y Claude & Innovaciones y figura de egresos & Alto & Medio & Alex Aravena revisó inversiones y costos con el modelo financiero y el T-19. Confirmó la diferencia entre beneficio para el cliente, ingreso variable y VAN base, además de revisar la figura mensual. \\
\end{tablalafrox}
